% mode: ini-etex
% Aimed at \csname and \ifcsname, whose characters from a token list are taken in bulk (§398, §1767), at how expand parts its
% commands between the backups of §388 and none, and at the conditionals of cond_other.
\input prelude
\errorcontextlines=1000
\catcode`\&=4
\catcode`\_=10
\scrollmode
\def\report#1{\immediate\write16{[#1]}}
\def\num#1{\report{\detokenize{#1} = \number\numexpr#1\relax}}
\def\relaxlike{\relax}
\let\letrelax=\relax
\let\amp=&
\def\five{5}
\def\oh{}
\font\nullf=nullfont

% ---- \csname and \ifcsname: characters from the list, a parameter, expansions, and what must not be taken in bulk
\def\c#1{\expandafter\meaning\csname #1\endcsname}
\report{\c{abc}}
\report{\c{a b}}
\report{\c{a1 b2}}
\def\namea{abc}
\report{\c{\namea}}
\report{\c{\namea\namea d}}
\report{\c{}}
\report{\c{x}}
\report{\c{\oh}}
\expandafter\def\csname abc\endcsname{ABC}
\report{\c{abc}}
\report{\c{\namea}}
\report{\c{\c{abc}}}
\report{\csname abc\endcsname}
\def\t#1{\csname a#1c\endcsname}
\report{\t{b}}
\def\t{\csname ab{c}\endcsname}
\report{\t}
\def\t{\csname ab\relax c\endcsname}
\report{\t}
\def\t{\csname ab\undefinedafter c\endcsname}
\report{\t}
\def\t{\csname ab}
\report{\t c\endcsname}
\def\t{\csname ab\ifnum 1<2 c\fi\endcsname}
\report{\t}
\def\t{\csname ab\ifcsname abc\endcsname c\fi\endcsname}
\report{\t}
\def\t{\ifcsname abc\endcsname yes\else no\fi}
\report{\t}
\def\t{\ifcsname ab\endcsname yes\else no\fi}
\report{\t}
\report{\ifcsname ab\endcsname yes\else no\fi}
\report{\ifcsname\namea\endcsname yes\else no\fi}
\report{\ifcsname\endcsname yes\else no\fi}
\report{\ifcsname x\endcsname yes\else no\fi}
\report{\ifcsname y\relax\endcsname yes\else no\fi}
\report{\ifcsname y\undefinedafter z\endcsname yes\else no\fi}
\tracingcommands=2
\report{\csname abc\endcsname}
\report{\ifnum 1<2 yes\fi}
\report{\ifx aa yes\fi}
\report{\ifcsname abc\endcsname yes\fi}
\tracingifs=1
\tracingcommands=0
\report{\csname abc\endcsname}
\report{\ifnum 1<2 yes\else no\fi}
\report{\unless\ifnum 1<2 yes\else no\fi}
\report{\ifcase 1 a\or b\else c\fi}
\report{\ifcase 3 a\or b\else c\fi}
\report{\ifcase 3 a\or b\fi}
\tracingifs=0
\def\t{\csname\expandafter\endcsname\undefinedafter}
\report{\t}
\tracingassigns=1
\report{\csname newafterassign\endcsname}
\tracingassigns=0

% ---- commands expand leaves to the procedure as WEB has it, and with the tracing that is tested first
\tracingcommands=2
\expandafter\relax\relax
\expandafter\expandafter\expandafter\relax\expandafter\relax\relax
\report{\ifnum\numexpr 1\relax<2 yes\else no\fi}
\tracingcommands=0
\tracingonline=1
\def\t{\the\numexpr 1+\number\numexpr 2+\romannumeral\numexpr 1\relax\relax\relax}
\report{\t}
\def\t{\number\dimexpr 1pt+\numexpr 2*3\relax sp\relax}
\report{\t}
\ifcase\numexpr 2\relax a\or b\or \report{c}\else d\fi
\count0=\numexpr 7\relax\report{\the\count0}
\count\numexpr 1\relax=\numexpr 8\relax\report{\the\count1}
\multiply\count0 by \numexpr 2\relax\report{\the\count0}
\end
